\section{Design and the Coupling Theorem}
\label{sec:design}

\subsection{System design}
\label{sec:design-overview}

NeuralPSI is a two-party protocol: a client-side plaintext Siamese CNN, a public fixed binariser to a
$d$-bit code $\xt \in \{0,1\}^d$ ($d=128$), and FLPSI (random $w$-subset sub-sampling, count-threshold
accept, label-only output) deciding whether $\xt$ matches an enrolled $\yt_i$ and revealing only
$\ell_i$; the functionality and accept model are due to~\cite{UzunEtAl21}, and \S\ref{sec:security}
proves what each corrupt party learns. This section is about what happens \emph{before} the protocol
runs: how the code is chosen, and why that choice governs both accuracy and communication through a
single scalar. \S\ref{sec:binarisers} gives two binarisers for the same target, derived next.

\subsection{Notation and setup}
\label{sec:coupling-setup}

For a probe/enrolled pair write $H = \lVert \xt \oplus \yt\rVert_1$ and $D_i = \mathbf{1}[\xt_i \neq
\yt_i]$, so $H = \sum_{i=1}^d D_i$. A single trial collides with probability $q(H) =
\binom{d-H}{w}/\binom{d}{w} \approx (1-H/d)^{w}$ (hypergeometric; the replacement approximation is tight
for $w\ll d$), so $K \mid H \sim \mathrm{Binomial}(T, q(H))$ and $\Pr[\text{accept}\mid H] = \Pr[K\ge
t]$. With $P_{\mathrm{gen}}, P_{\mathrm{imp}}$ the genuine and impostor laws of $H$, FAR and FRR are the
exact functionals $\sum_H P_{\mathrm{imp}}(H)\Pr[\text{accept}\mid H]$ and $1-\sum_H
P_{\mathrm{gen}}(H)\Pr[\text{accept}\mid H]$; write $q_g = \mathbb{E}_{P_{\mathrm{gen}}}[q(H)]$ and
$q_i = \mathbb{E}_{P_{\mathrm{imp}}}[q(H)]$ for the aggregated per-trial accept masses. Four assumptions
underlie what follows. (A1) Acceptance depends
on a pair only through $H$, since $S_k$ is uniform (exact by construction). (A2) Given $H$, a
sub-sample collides with probability $q(H)$ despite bit dependence: measured $\hat q(H)/q(H)$ has median
$1.002$ over $H \in [5,40]$ at $w=14$, so dependence shows up in the distribution of $H$, not in
$\Pr[\text{accept}\mid H]$, isolating the modelling problem to $P_{\mathrm{imp}},P_{\mathrm{gen}}$. (A3)
The masks are independent, so $K \mid H \sim \mathrm{Binomial}(T, q(H))$. (A4) Online per-query
communication scales linearly in $T$.

\subsection{The coupling analysis}
\label{sec:coupling}

A single scalar functional of the code, the variance of the impostor Hamming law, governs the
achievable (FAR, FRR) operating point and the number of sub-sampling rounds $T$ needed to reach it,
and hence, by (A4), the online communication. A code that minimises this functional is accuracy- and
communication-optimal at once. We state the results here; full proofs are in the extended version.

\begin{lemma}[Cost is a functional of the impostor-Hamming law]
\label{lem:jensen}
Let $g(h) = (1-h/d)^w$ on $h \in [0,d]$, $w \ge 2$. Then $g$ is convex and strictly decreasing, and
the per-trial impostor accept mass satisfies
$q_i = \mathbb{E}_{P_{\mathrm{imp}}}[g(H)] \ge g(\mathbb{E}[H]) = (1-\mathbb{E}[H]/d)^{w}$. The gap
$q_i - g(\mathbb{E}[H])$ is monotone under mean-preserving spreads of $P_{\mathrm{imp}}$ (the convex
order), and to second order equals $\tfrac12 g''(\mathbb{E}[H])\,\mathrm{Var}(H)$ plus a third-moment
term. At fixed mean $\mathbb{E}[H]=d/2$, larger impostor variance means larger $q_i$. ($g$ is decreasing and
convex for $w\ge2$, so Jensen and the convex order apply, with the variance term the second-order
Taylor coefficient.)
\end{lemma}

The rigorous statement is the convex-order one; plain variance is its second-order, operational
surrogate, and is what we track empirically. The paper's central diagnostic, an impostor law with
$\sigma = 20.11$ against the uniform ideal $5.66$ at the same mean $64$, is exactly such a
mean-preserving spread, which is why it inflates $q_i$ and, by Proposition~\ref{prop:tmin}, the
required $T$.

\begin{proposition}[Neyman--Pearson optimality of the count threshold]
\label{prop:np}
Fix $(w,T)$ and let $K$ be the colliding-trial count for a pair with per-trial mass
$q \in \{q_i, q_g\}$, $q_g > q_i$. $K$ is a sufficient statistic for deciding genuine versus impostor
within $K \sim \mathrm{Binomial}(T,q)$, and the likelihood ratio
$\Lambda(k) = (q_g(1-q_i)/q_i(1-q_g))^{k}\,((1-q_g)/(1-q_i))^{T}$ is monotone increasing in $k$. By
Neyman--Pearson with a monotone likelihood ratio, ``accept iff $K \ge t$'' is uniformly most powerful
at its size: thresholding $K$ minimises FRR at any given FAR.
\end{proposition}

So FLPSI's count-threshold rule is the optimal detector for the statistic the protocol computes, not
an ad hoc heuristic. It does not by itself say how many rounds $T$ are needed; that is next.

\begin{proposition}[Required $T$]
\label{prop:tmin}
Model $K \mid \text{impostor} \sim \mathcal{N}(Tq_i, Tq_i(1-q_i))$ and $K \mid \text{genuine} \sim
\mathcal{N}(Tq_g, Tq_g(1-q_g))$. Placing $t$ to hit $\mathrm{FAR}=\alpha$ and $\mathrm{FRR}=\beta$
requires, with $z_\alpha=\Phi^{-1}(1-\alpha)$, $z_\beta=\Phi^{-1}(1-\beta)$,
\begin{equation}
\label{eq:tmin-gauss}
T_{\min} \;\approx\; \frac{\bigl[z_\alpha\sqrt{q_i(1-q_i)} + z_\beta\sqrt{q_g(1-q_g)}\bigr]^2}
{(q_g - q_i)^2}.
\end{equation}
\end{proposition}

$T_{\min}$ scales as $1/(q_g-q_i)^2$, so shrinking $q_i$, which Lemma~\ref{lem:jensen} ties to the
impostor variance, widens the gap and is the dominant lever on the rounds the protocol needs.

\paragraph{Finite-$T$ caveat.} Proposition~\ref{prop:tmin} is asymptotic. At the small $T$ deployed
($T=18$) we use the exact binomial tail for design and validate against the real cryptographic binary:
at $(w,t,T) = (6,3,18)$ the feature-head code (\S\ref{sec:binarisers}) gives measured TAR $97.0\%$ and
measured FAR $8.02\times10^{-3}$ against a predicted $8.12\times10^{-3}$, and earlier codes matched
their binomial predictions essentially exactly. We present the asymptotics as design guidance and the
executed protocol as ground truth (\S\ref{sec:eval}).

\subsubsection{The variance decomposition and the covariance floor}

Lemma~\ref{lem:jensen} and Proposition~\ref{prop:tmin} reduce the design problem to one target:
minimise $\mathrm{Var}(H_{\mathrm{imp}})$ at a mean that keeps $q_g$ usable. The next theorem shows
where that variance comes from. For an impostor pair write $p_i = \Pr[D_i=1]$, so $\mathbb{E}[H] =
\sum_i p_i$ and $\mathrm{Var}(D_i) = p_i(1-p_i)$.

\begin{theorem}[Variance decomposition and the covariance floor]
\label{thm:variance}
\begin{equation}
\label{eq:vardecomp}
\mathrm{Var}(H_{\mathrm{imp}}) \;=\; \sum_{i=1}^{d} p_i(1-p_i) \;+\; \sum_{i \neq j}
\mathrm{Cov}(D_i,D_j).
\end{equation}
The diagonal term obeys $p_i(1-p_i)\le 1/4$ with equality iff $p_i=1/2$, so $\sum_i p_i(1-p_i)\le d/4$,
attained at full balance. The uniform/max-entropy code has every $p_i=1/2$ and every
$\mathrm{Cov}(D_i,D_j)=0$, giving $H \sim \mathrm{Binomial}(d,1/2)$, $\mathrm{Var}=d/4$,
$\sigma=\sqrt{d}/2$ ($=5.657$ at $d=128$). Every unit of variance above the floor $d/4$ is therefore
positive inter-bit covariance. (Bilinearity of variance on $H=\sum_i D_i$ gives \eqref{eq:vardecomp};
$p(1-p)\le1/4$ is maximised at $p=1/2$; the floor identity is \eqref{eq:vardecomp} rearranged.)
\end{theorem}

The three codes instantiate this against the balanced diagonal ceiling $d/4=32$. Super-Bit is balanced
(mean $64.05$, $\sigma=20.11$, $\mathrm{Var}=404.5$), so $\sum_i \mathrm{Var}(D_i) \approx 32$ and the
residual covariance $\sum_{i\neq j}\mathrm{Cov}(D_i,D_j) \approx 372.5 \gg 0$. The ortho-thermometer
code (mean $47.27$, $\sigma=14.60$) is not fully balanced, so its diagonal ceiling is already below
$32$; the feature-head code (mean $64.09$, $\sigma=6.72$, $\mathrm{Var}=45.2$) is balanced and near the
floor, with covariance surplus $\approx 13.2$.

\begin{corollary}[What to minimise]
\label{cor:target}
Among impostor laws with mean $d/2$, $\mathrm{Var}(H_{\mathrm{imp}})$, and hence by
Lemma~\ref{lem:jensen} $q_i$, by Proposition~\ref{prop:tmin} $T$, by (A4) online communication, and
hence the achievable $(\mathrm{FAR},\mathrm{FRR})$, is minimised by the balanced and decorrelated code
whose impostor law is $\mathrm{Binomial}(d,1/2)$ with $\sigma = 5.66$. A single improvement to the code
improves accuracy and communication together. The measured impostor $\sigma$ ladder
$20.11 \to 14.60 \to 6.72$ tracks the EER ladder $1.89\% \to 0.87\% \to 0.17\%$ and the per-query
communication at $N=5000$, $4317 \to 2377$\,KB ($-45\%$).
\end{corollary}

\subsubsection{The feature-head objective relaxes Corollary~\ref{cor:target}}

Corollary~\ref{cor:target}'s population program (minimise $\sum_{i\neq j}\mathrm{Cov}(D_i,D_j)$ subject
to $\mathbb{E}[D_i]=1/2$ and small $\mathbb{E}_{\mathrm{gen}}[H]$) is realised by training the
feature-head against $\mathcal{L} = \lambda_{\mathrm{gen}}\mathcal{L}_{\mathrm{gen}} +
\lambda_{\mathrm{bal}}\mathcal{L}_{\mathrm{bal}} + \lambda_{\mathrm{dec}}\mathcal{L}_{\mathrm{dec}} +
\lambda_{q}\mathcal{L}_{\mathrm{STE}}$ with $(\lambda_{\mathrm{dec}},\lambda_{\mathrm{bal}},\lambda_q) =
(30,5,0.1)$.

\begin{theorem}[The feature-head loss is a sampled relaxation of the population program]
\label{thm:relaxation}
Term by term, $\mathcal{L}$ is a differentiable, sample-average surrogate of that program:
$\mathcal{L}_{\mathrm{bal}}$ pins $\mathbb{E}[D_i]=1/2$; $\mathcal{L}_{\mathrm{dec}}$ penalises the
off-diagonal bit covariance, driving $\sum_{i\neq j}\mathrm{Cov}(D_i,D_j)\to 0$;
$\mathcal{L}_{\mathrm{gen}}$ keeps $\mathbb{E}_{\mathrm{gen}}[H]$ small; and the straight-through
estimator makes the sign binarisation trainable.
\end{theorem}

Because $\mathrm{Var}(H_{\mathrm{imp}})$ depends only on $\{p_i\},\{\mathrm{Cov}(D_i,D_j)\}$
(Theorem~\ref{thm:variance}), balance plus decorrelation is a tight, convex surrogate for it. Two
caveats hold: the map from weights to those moments is non-convex and the straight-through estimator
biased, so the convexity is of the objective as a functional of the bit distribution, not of the
optimisation landscape (no global-optimality claim); and a variance floor is not a law, so zeroing
linear covariance does not force the impostor law to $\mathrm{Binomial}(d,1/2)$, read instead as
``moves toward the uniform ideal'', consistent with $\sigma \to 6.72 \approx 5.66$. Full proof in the
extended version.

The theorem co-designs the code \emph{for} the protocol's own optimal decision rule
(Proposition~\ref{prop:np}). One tension it does not resolve: tightening $\sigma$ can raise the
\emph{genuine} Hamming distance, which lowers $q_g$, so the single-finger EER and communication wins
are robust while multi-sample fusion stays a tunable trade-off.

\subsubsection{Effective bits and the collision floor}

The variance functional also gives a degrees-of-freedom estimate: matching the variance of the
impostor Hamming fraction $H/d$ to a binomial gives the Daugman effective-bits
estimator~\cite{Daugman03}
\begin{equation}
\label{eq:meff}
m_{\mathrm{eff}} = \frac{p(1-p)\,d^2}{\sigma^2} = \frac{d^2}{4\sigma^2} \quad\text{at balance},
\end{equation}
with $m_{\mathrm{eff}}=d$ for the ideal code. A 16-dimensional embedding bottleneck caps Super-Bit at
$\approx10$ effective bits of 128, while the wider convolutional representation the feature-head codes
lifts it to $90.7$; numeric values with a participation-ratio check are in \S\ref{sec:effbits-params}.

\begin{corollary}[Collision floor]
\label{cor:floor}
Since $q(0)=1$, $\Pr[\text{accept}\mid H=0]=1$ for every $(w,t,T)$, so the zero-Hamming impostor
collision probability $p_0$ lower-bounds the per-record FAR for any parameterization and propagates to
$\mathrm{FAR}_{\mathrm{query}}=1-(1-p_0)^N$. Bounding $p_0$ away from zero is a coding question governed
by the \emph{effective} dimension $m_{\mathrm{eff}}$, not the nominal $d$ (sphere-packing/Singleton/
Plotkin on the effective code), so raising $d$ without raising $m_{\mathrm{eff}}$ leaves the floor
essentially unchanged, as a Super-Bit dimension sweep shows.
\end{corollary}

As a diagnostic, Super-Bit has $2/1{,}438{,}800$ zero-Hamming impostor collisions on the held-out
probe$\times$gallery set ($p_0 = 1.39\times10^{-6}$), while both co-designed codes drive the count to
$0$; we recommend reporting the held-out $H=0$ count for any binariser.

\subsection{Two binariser instantiations}
\label{sec:binarisers}

The target is one thing: balance the bits and decorrelate them, at a mean that keeps genuine Hamming
distance small. We give two binarisers at different points on the training-cost/accuracy trade-off.

\paragraph{Training-free: the metric-aware ortho-thermometer.}\label{sec:ortho-thermo}
This code needs no training beyond the extractor in use. It replaces a single random-hyperplane sign
per output dimension with a per-dimension thermometer code from block-orthonormalised projection
directions and population calibration statistics, so Hamming distance under the accept rule tracks the
extractor's native angular (cosine) metric rather than a raw per-axis threshold. Being training-free it
is metric-\emph{aware}, not metric-\emph{learned}: it corrects the mismatch between an angular embedding
and a Hamming matcher without touching the network, dropping the impostor standard deviation from
$20.11$ to $14.60$ by improving balance and, less so, decorrelation.

\paragraph{Quantization-aware: the learned feature-head.}\label{sec:featurehead}

The feature-head instantiates Theorem~\ref{thm:relaxation} directly on the wider convolutional
representation rather than the 16-dimensional bottleneck feeding Super-Bit. A bit-balance loss pins
each bit's mean to $1/2$ (measured balance $\to 0.499$); a bit-decorrelation loss penalises the
off-diagonal bit covariance, targeting the $\sum_{i\neq j}\mathrm{Cov}(D_i,D_j)$ term of
Theorem~\ref{thm:variance}; a genuine-closeness loss keeps $\mathbb{E}_{\mathrm{gen}}[H]$ small; and a
straight-through estimator gives a trainable gradient through the sign binarisation. The result is the
$\sigma = 6.72$, $m_{\mathrm{eff}} \approx 91$ code used throughout \S\ref{sec:eval}, the closest of
the three to the balanced, decorrelated target, with the caveats of Theorem~\ref{thm:relaxation}
carried forward.
